% ======================================================================
\section{Feedback from the molecules on MOCASSIN}
\label{sec:feedback}
% ======================================================================

The molecular content computed in Section~\ref{sec:chemistry} acts back on
the radiation field, the ionisation and the temperature, which closes the
loop of Figure~\ref{fig:chain}. It does so through three channels only: the
fraction of each element held in molecules, the \Htwo{} photoionisation
continuum, and the heating and cooling terms. No molecular line enters the MC
transfer.

\subsection{Elements locked in molecules}

The whole handback for an element is the fraction of it \emph{not} in
molecules (\rtn{setElemAvail}, \src{pdr\_mod.f90:\allowbreak{}5120--5144}, called by
\rtn{pdrToMocassin}, \src{:\allowbreak{}5038--5104}):
\begin{equation}
  \rtn{elemAvail}_X = \min\!\Bigl[\max\Bigl(1-\frac{M_X}{A_X},\,0\Bigr),\,1\Bigr],
  \qquad M_X = \sum_i \nu_{X,i}\,x_i ,
\end{equation}
for the five linked elements H, He, C, O and Mg, with $A_X$ MOCASSIN's own
abundance ($A_{\mathrm H}=1$) and $\nu_{X,i}$ the number of atoms of $X$ in
species $i$, read from the species' formula for every species other than $X$
and $X^+$ themselves (\src{:\allowbreak{}585--603}, \dlog{D-402}). MOCASSIN's
\rtn{ionBalance} constrains the sum of the element's ion stages to it
(\src{update\_mod.f90:\allowbreak{}3130--3134}): carbon in CO therefore appears in no
atomic stage, emits no atomic line and gives no free electron. The chemistry
writes no ionisation state and no temperature (\dlog{D-193}). After each
iteration \rtn{rescaleIonToAvail} (\src{pdr\_mod.f90:\allowbreak{}2764--2803}) rescales the
stages of every element to \rtn{elemAvail}, keeping MOCASSIN's ratios between
them, so that the hydrogen budget
$x(\mathrm H^0)+x(\mathrm H^+)+2x(\mathrm{H_2})=1$ holds in every cell (up to
the hydrogen in other molecules) before the next transfer builds its opacity
(\dlog{D-220}). It prints the size of the pre-rescale defect (\texttt{[D220]}),
which is the informative test of the budget: since the kept second solve
(below) no cell of either fixture is off by more than 1\% in the written model
(\dlog{L-29}).

The network channels classified in Section~\ref{sec:chemistry} (sources and
sinks of $X$ and $X^+$ to and from molecules, and net conversions between
them) enter \rtn{ionBalance}'s rate matrix as sinks on the diagonal and
sources on the right-hand side, the sum still constrained to \rtn{elemAvail}
(\src{update\_mod.f90:\allowbreak{}3255--3300}, \dlog{D-402}). They are those of the
previous chemistry, i.e.\ one iteration behind (\dlog{L-47}).
Cosmic-ray ionisation, which MOCASSIN otherwise lacks, is added as MOCASSIN's
own physics with the network's coefficients: $\alpha\,\zeta/\zeta_0$ for direct
ionisation and
$\alpha\,(\zeta/\zeta_0)(T/300\,\mathrm K)^\beta\gamma/(1-\omega)$ for the
cosmic-ray-induced photons, $\zeta_0=1.3\times10^{-17}$\,s$^{-1}$, for every
linked neutral (\src{update\_mod.f90:\allowbreak{}2916--2930}; \dlog{D-193}, \dlog{D-280}).

\subsection{\texorpdfstring{\Htwo}{H2} in the gas opacity}

$n(\mathrm{H_2})$ (\rtn{H2den}) adds the \Htwo{} photoionisation continuum
above 15.4\,eV to the gas opacity of the next transfer
(\rtn{addOpacity}, \src{ionization\_mod.f90:\allowbreak{}435--472}, \dlog{D-107}),
\begin{equation}
  \kappa_\nu \mathrel{+}= n(\mathrm{H_2})\,\sigma_{\mathrm{H_2}}(\nu),\qquad h\nu\ge15.4\unit{eV},
\end{equation}
with the cross section fitted by \citet{Yan1998}, as revised in their erratum
\citep{Yan2001} (\rtn{yanH2xSec},
\src{ionization\_mod.f90:\allowbreak{}385--411}; zero below 15.4\,eV, as in Cloudy, and
clipped at zero just above threshold where the fit goes negative). The
comment above the opacity loop (\src{:\allowbreak{}444--453}) still describes ``a
single power law anchored at threshold''; the function actually used is the
piecewise fit of \citet{Yan1998}. This continuum matters only for
hydrogen-ionising photons that reach molecular gas. No \Htwo{}, CO or
H\,\textsc{i} Lyman-series \emph{line} opacity is added to the MC transfer:
line absorption is represented only through the shielding columns.

\subsection{Heating and cooling}
\label{sec:thermal}

MOCASSIN solves the thermal balance of every cell, heating equal to cooling
per H nucleus, with \rtn{thermBalance} (\src{update\_mod.f90:\allowbreak{}1734--2456})
inside its temperature search \rtn{iterateT} (\src{:\allowbreak{}609--1191}). Its own
terms are the photoionisation heating by every ion on the MC field; the grain
photoelectric heating (\rtn{setPhotoelHeatCool},
\src{update\_mod.f90:\allowbreak{}1460--1556}), summed over every grain species and
size from the MC field above the photoelectric threshold (work function plus
grain potential), with the yield of \citet{WeingartnerDraine2001}
(\dlog{D-289}) and the grain charge of MOCASSIN's own charge balance, whose
electron sticking was set to the same paper's value (\dlog{D-326}); the
recombination, free--free and collisionally excited line cooling of every ion;
and the gas--grain collisional exchange (\rtn{setDustGasCollHeatCool},
\src{:\allowbreak{}1562--1731}), which omits collisions with \Htwo{} (\dlog{L-9}).
The molecules enter it in three ways.

\begin{enumerate}
\item \textbf{Coolants.} The coolants named by \kw{pdrCoolants} (in the
  fixtures \Cp, C, O and $^{12}$CO, LAMDA-format files, \citealt{Schoier2005})
  are treated in two ways.
  \begin{itemize}
  \item \Cp, C and O are cooled by MOCASSIN's own model atoms (its electron
    collision data and all its levels), to which the coolant files' rates for
    collisions with H, \Hp, He and \Htwo{} and the chemistry's escape
    probability of each transition are added (\src{emission\_mod.f90:\allowbreak{}1861--1887};
    \src{common\_mod.f90:\allowbreak{}1489--1635}; \dlog{D-363}). They are solved at every
    trial temperature, and the same atoms give the reported line fluxes, so
    the thermal balance and the lines share one atom per ion.
  \item CO, which MOCASSIN has no atom for, is solved by 3D-PDR's
    escape-probability level-population solver in \rtn{pdrStep} at 16
    temperatures log-spaced over [\kw{pdrTmin}, \kw{pdrTmax}], a factor 1.94
    apart for the defaults (\rtn{initPDRcoolTab},
    \src{pdr\_mod.f90:\allowbreak{}2713--2732}, \dlog{D-215}). The temperature search reads
    this table by log--log interpolation, clamped at its ends (\rtn{pdrCoolAt},
    \src{common\_mod.f90:\allowbreak{}1445--1486}), and adds it to MOCASSIN's cooling
    (\src{update\_mod.f90:\allowbreak{}2232--2249}). For CO the table misses an exact solve
    at the cell's temperature by a mean $|\Delta\ln\Lambda|$ of 0.02
    (\texttt{hiipdr}) and 0.004 (\texttt{pn150k}) per cell (\dlog{L-36}).
  \end{itemize}
  In both the level populations are in statistical equilibrium with the
  colliders H, \Hp, $e^-$, He and \Htwo{} at the chemistry's abundances (the
  linked ions and $n_e$ being MOCASSIN's). The escape probability of a
  transition $u\to l$ is the mean over the twelve line rays of
  Section~\ref{sec:lineRays} of $\beta=(1-e^{-\tau})/\tau$, with
  \begin{equation}
    \tau = \frac{A_{ul}\,c^3}{8\pi\nu_{ul}^3\,b}\sum_{\mathrm{cells}}
           \Bigl(n_l\frac{g_u}{g_l}-n_u\Bigr)\ell ,
  \end{equation}
  each crossed cell contributing its own level populations as last stored
  (\src{pdr/\allowbreak{}escape\_probability.F90:\allowbreak{}121--137, 196--202}; \dlog{D-377}). With
  \kw{-DSOBOLEV} the rate equations use the equivalent net-rate form, the
  spontaneous rate $A_{ul}\beta$ and the radiative rates driven by
  $\beta J_{\mathrm{bg}}$ (\src{:\allowbreak{}219--233}), where the background
  $J_{\mathrm{bg}}$ is MOCASSIN's own $J_\nu$ at the line plus the CMB
  (\dlog{D-377}). The cooling is the net escaping power,
  $\Lambda_{ul}=A_{ul}h\nu_{ul}n_u\beta\,(S_{ul}-J_{\mathrm{bg}})/S_{ul}$ with
  $S_{ul}$ the line source function (\src{:\allowbreak{}209--211}).
\item \textbf{Extra heating.} 3D-PDR's heating routine \rtn{CALC\_HEATING}
  (\src{pdr/\allowbreak{}heatingfunctions.F90}) is called at \src{pdr\_mod.f90:\allowbreak{}4499}
  with the local field $\chi$ at MOCASSIN's $T_e$. The terms MOCASSIN computes
  itself are excluded: grain and PAH photoelectric heating (3D-PDR's
  \citealt{BakesTielens1994} expression with the modifications of
  \citealt{Wolfire2003}, and its alternative \citealt{WeingartnerDraine2001}
  fit), carbon photoionisation heating and gas--grain collisions; these are
  the only terms that read $\chi$. The rest, \rtn{pdrHeatXtra}
  (\src{pdr\_mod.f90:\allowbreak{}4545--4546}), is added to MOCASSIN's heating
  (\src{update\_mod.f90:\allowbreak{}2391}) and held fixed during the temperature search:
  \begin{itemize}
  \item \emph{\Htwo{} terms from the level model} (\src{pdr\_mod.f90:\allowbreak{}4504--4528},
    \dlog{H-013}): formation, $R_{\mathrm f}\,n_{\mathrm H}\,n(\mathrm H)\,E_{\mathrm{kin}}$,
    with $E_{\mathrm{kin}}$ the kinetic energy of a newly formed molecule from the
    formation distribution; dissociation, $n(\mathrm{H_2})\sum_l n_l H^{\mathrm{dis}}_l$,
    the kinetic energy of the atoms; and the net collisional term
    $n(\mathrm{H_2})\,\mathcal{H}_{\mathrm{coll}}$, which is the heating by the
    de-excitation of pumped and newly formed molecules and, where the gas is
    warmer than the level populations, \Htwo's line cooling
    (Section~\ref{sec:solvecell}). They replace 3D-PDR's fixed energies
    (1.5\,eV per formation, which the vendored code attributes to
    \citealt{HollenbachTielens1999}; 0.4\,eV per dissociation;
    $9\times2.2$\,eV per dissociation for pumping, with the critical-density
    factor it attributes to \citealt{HollenbachMcKee1979}), which are kept
    beside them for comparison (\src{h2Cells.out});
  \item cosmic-ray heating, $\Gamma_{\mathrm{CR}}=9.4\,\mathrm{eV}\times\zeta\,n(\mathrm{H_2})$,
    i.e.\ $1.5\times10^{-11}\zeta\,n(\mathrm{H_2})$\,erg\,cm$^{-3}$\,s$^{-1}$, the
    expression of \citet[their eq.~A23]{TielensHollenbach1985}: about 8\,eV of
    heat per primary ionisation in molecular gas, helium included, with $\zeta$
    the rate per \Htwo{}. Ionisations in atomic gas deposit no heat in this
    term, which \citet{TielensHollenbach1985} justify by cosmic-ray heating
    being unimportant there;
  \item turbulent dissipation, $\Gamma_{\mathrm{turb}}=3.5\times10^{-28}\,(v_{\mathrm{turb}}/\mathrm{km\,s^{-1}})^3\,
    (L_{\mathrm{turb}}/\mathrm{pc})^{-1}\,n_{\mathrm H}$\,erg\,cm$^{-3}$\,s$^{-1}$ with
    $L_{\mathrm{turb}}=5$\,pc fixed and $v_{\mathrm{turb}}$ = \kw{pdrVturb}, which the
    vendored code attributes to \citet{Black1987}: $7\times10^{-29}$\,erg\,s$^{-1}$
    per H nucleus at the default 1\,km\,s$^{-1}$;
  \item the energy released by exothermic reactions (for example the
    dissociative recombination of HCO$^+$, H$_3^+$ and H$_3$O$^+$, and He$^+$ with
    \Htwo{} and CO),
    the reactions being found automatically in the network, after
    \citet{Clavel1978}, who found that such reactions could contribute
    significantly to the heating if their energy is released entirely as
    kinetic energy, as is assumed here.
  \end{itemize}
\item \textbf{Photoelectric heating} is MOCASSIN's own, from its grains, as
  described above (\dlog{D-289}, \dlog{D-326}); 3D-PDR's PAH term is not used
  (\dlog{D-134}, \dlog{D-352}). The model grains include no PAHs (\dlog{L-7}),
  so the heating lacks the contribution of PAHs and very small grains that
  dominates in the \citet{BakesTielens1994} prescription of most PDR codes
  and of the code comparison \citep{Rollig2007}; the Cloudy reference models
  omit them too. Against Cloudy the photoelectric efficiency per H per unit
  field is 0.66--1.10 with no trend with depth on \texttt{pn150k}
  (\dlog{D-304}); on \texttt{hiipdr} the gas temperature at $\AV=0.73$ and 0.97
  is 1.07 and 1.01 times Cloudy's since the sticking fix, the deeper gas
  remaining colder (\dlog{L-22}).
\end{enumerate}

The gas temperature found this way is the one the next chemistry runs at,
and, one iteration later, the level model. It acts on the molecules through
the rate coefficients, through the \Htwo{} formation coefficient (the thermal
speed of the H atoms and their sticking probability,
Section~\ref{sec:formation}), through the collisional excitation of \Htwo{}
and of the coolants, and through the CO cooling it controls. The dust
temperature acts through the formation efficiency (Section~\ref{sec:Td}).
These are channels by which the field sets the molecular content even where
no FUV photon reaches the molecules.

\subsection{The second ionisation solve}

After the chemistry, \rtn{ionBalance} is solved a second time in every cell,
at the temperature of the first solve, on the channels and the molecular
content this iteration's chemistry has just left; it is counted against the
first and the model is put back (\rtn{ionSecondSolve},
\src{iteration\_mod.f90:\allowbreak{}3361--3434}, called at \src{:\allowbreak{}2168}; \dlog{D-404}). A
control pass on the old terms, before the chemistry (\src{:\allowbreak{}1274}), moves
nothing. On the finished model the second solve's stages replace the first's
(\rtn{ionKeepSecond}, \src{:\allowbreak{}3455--3521}), so the written ions are consistent
with the last chemistry, and the level model's final call runs on them
(\dlog{H-029}). The second solve removes the one-iteration delay of
\dlog{L-47}: the neutral cells whose solver $X^+/X$ is more than 10\% from
MOCASSIN's fall from 28.4\% to 0.0\% for hydrogen on \texttt{hiipdr}, and to
0.0\% for every linked element except oxygen there (9.7\%, \dlog{L-50}). Neither
the temperature nor the chemistry is solved again on the second solve's ions,
and in the iterations before the last the transfer, the thermal balance and
the chemistry still run on the first solve's.
